\section*{Broader impact}

\subsection{Intended use}

This work is a research contribution on cost-sensitive verification,
recovery, and process supervision for scientific agents. The controller and
training framework are designed for research evaluation under controlled
fault injection and independent adjudication. They are not a deployed safety
system and should not be treated as one.

\subsection{Beneficiaries}

Primary beneficiaries are researchers developing reliable agent systems who
require methods that (a)~separate belief from decision, (b)~account for
verification and recovery costs, (c)~evaluate under frozen splits with
independent adjudication, and (d)~test for clean-task regression. The
verified-intervention dataset may also benefit the broader community
studying reinforcement learning from process-level feedback.

\subsection{Misuse and dual-use risks}

A cost-sensitive verification policy could be used to minimize the
appearance of failure (by selectively verifying only cases likely to pass)
rather than to maximize scientific validity. The requirement that the
primary endpoint is independently adjudicated scientific success, not
self-reported success, mitigates this risk. The framework could also be
adapted to reduce the cost of deploying unreliable agents in high-stakes
settings without improving their fundamental reliability; the clean-task
regression gate and the requirement for independent adjudication are
necessary but not sufficient protections.

\subsection{Cost-model distributional effects}

The cost weights ($\lambda_v$, $\lambda_{\text{compute}}$,
$\lambda_{\text{human}}$, $\lambda_{\text{damage}}$) encode value judgments
about the relative importance of verification effort, compute, human time,
and irreversible damage. Different stakeholders may reasonably choose
different weights. The paper reports sensitivity to these choices rather
than prescribing one set of weights. Users who adopt the controller with
their own cost weights should re-evaluate the primary results under those
weights.

\subsection{Privacy and security controls}

The framework does not collect personal data. Tasks are drawn from public
scientific repositories. The verified-intervention dataset records
structured transition metadata; raw scientific data referenced through
content hashes is governed by the R1 benchmark's access and licensing
policies.

\subsection{Mitigations and residual risks}

Structural mitigations include: separation of training and evaluation
verifiers, the clean-task non-inferiority gate, the verifier-substitution
audit, the reward-channel audit, frozen splits and deduplication controls,
and retention of ambiguous and negative outcomes. Residual risks include
the possibility that a model learns a recovery strategy that transfers to a
deployment domain with different failure modes, or that users ignore the
cost-sensitivity results and deploy the controller with an untested cost
configuration.
